%!TEX root = ./main.tex

\section[Example 1: DAS]{Example 1: Distributed Antenna System}

% TALK MAPPING: slides 9--12, transcript 10:28--16:00. This is the only application
% the speaker walked in full, and it is the right one for a first exposure: the idea
% is physical and obvious, the cost problem is a dollar figure, and the fix uses only
% A2 (replication) on the downlink and A3+A4 (cache, combine) on the uplink.

% Teaching note (120 s): draw the attenuation argument by hand before advancing --
% one antenna on the roof, signal dies on the way down, ground floor gets the worst
% throughput. Only then show the DAS building. Students should reach "so put the
% antennas where the users are" before the slide says it.
\begin{frame}[t]{Example 1: Distributed antenna system (DAS)}
  \vspace{0.1em}
  \centering
  \only<2->{{\small\headline{Replace the high-power antenna with low-power
    distributed ones}\par}}
  \only<1>{{\small\phantom{Replace the high-power antenna with low-power}\par}}

  \vspace{0.3em}
  \begin{tikzpicture}[scale=0.86,transform shape]
    % ---- Regular transmission: one rooftop antenna, coverage fades downwards.
    \node[font=\scriptsize\bfseries] at (1.0,3.8) {Regular transmission};
    \fill[softgreen,opacity=0.10] (1.0,3.1) circle (2.05);
    \fill[softgreen,opacity=0.12] (1.0,3.1) circle (1.40);
    \fill[softgreen,opacity=0.16] (1.0,3.1) circle (0.85);
    \draw[draw=ranblue!55,fill=blue!4,fill opacity=0.45] (0.2,0.0) rectangle (1.8,3.1);
    \foreach \y in {0.78,1.55,2.32}{\draw[ranblue!35] (0.2,\y) -- (1.8,\y);}
    \draw[softgray,thick] (0.85,3.1) -- (1.0,3.55) -- (1.15,3.1);
    \node[draw=ranblue,fill=ranblue,text=white,minimum width=0.6cm,
          minimum height=0.3cm,inner sep=1pt,font=\tiny\bfseries] at (0.62,0.25) {DU};

    % ---- DAS: many low-power antennas, uniform coverage.
    \only<2->{
      \node[font=\scriptsize\bfseries] at (4.3,3.8) {DAS};
      \foreach \y in {1.00,1.77,2.54}{
        \fill[softgreen,opacity=0.22] (3.9,\y) circle (0.72);
        \fill[softgreen,opacity=0.22] (4.9,\y) circle (0.72);}
      \draw[draw=ranblue!55,fill=blue!4,fill opacity=0.45] (3.5,0.0) rectangle (5.1,3.1);
      \foreach \y in {0.78,1.55,2.32}{\draw[ranblue!35] (3.5,\y) -- (5.1,\y);}
      \foreach \y in {0.95,1.72,2.49}{
        \foreach \x in {3.9,4.9}{
          \draw[softgray] (\x-0.1,\y+0.18) -- (\x,\y+0.42) -- (\x+0.1,\y+0.18);}}
      \node[draw=ranblue,fill=ranblue,text=white,minimum width=0.55cm,
            minimum height=0.3cm,inner sep=1pt,font=\tiny\bfseries] at (3.85,0.25) {DU};
      \node[draw=softgreen,fill=softgreen,text=white,minimum width=0.55cm,
            minimum height=0.3cm,inner sep=1pt,font=\tiny\bfseries] at (4.72,0.25) {DAS};
    }

    % ---- The market says this is not a niche.
    \only<2->{
      \begin{scope}[shift={(6.6,0.35)}]
        \draw[softgray] (0,0) -- (4.5,0);
        \foreach \i/\v/\y in {0/19.2/0.93, 1/23.5/1.14, 2/29.9/1.45,
                              3/38.1/1.85, 4/41.2/2.00}{
          \draw[draw=ranblue,fill=ranblue!55] (\i*0.86+0.22,0)
            rectangle (\i*0.86+0.72,\y);
          \node[font=\tiny,text=softgray] at (\i*0.86+0.47,\y+0.16) {\v};
        }
        \foreach \i/\yr in {0/2016,1/2017,2/2018,3/2019,4/2020}{
          \node[font=\tiny,text=softgray] at (\i*0.86+0.47,-0.22) {\yr};}
        \node[font=\tiny\bfseries,align=center] at (2.25,2.85)
          {DAS nodes operated by Boingo\\Wireless worldwide (thousands)};
      \end{scope}}
  \end{tikzpicture}

  \vspace{0.2em}
  \raggedright
  \only<2->{%
    {\small\textcolor{softgreen}{\textbf{Pros:}} Better, uniform coverage ---
      standard practice in enterprises, stadiums, airports\par}
    {\small\textcolor{talkred}{\textbf{Cons:}} Very high cost and deployment
      complexity --- hundreds of thousands to millions of dollars\par}}

  \source{https://www.statista.com/statistics/1041873/boingo-wireless-das-nodes/}{DAS node counts: Boingo Wireless, as cited in the authors' talk}
\end{frame}

% Teaching note (150 s): this is the payoff frame. Walk the downlink first (one packet
% in, two identical packets out = A2) and then the uplink (two packets in, hold one
% until the other arrives = A3, add them = A4, one packet out). Then read the line
% counts aloud. The point of the paper is that this is small.
\begin{frame}{The DAS middlebox: replicate down, combine up}
  \vspace{0.5em}
  \centering
  \begin{tikzpicture}[scale=0.95,transform shape]
    \node[dunode,minimum width=1.5cm,minimum height=0.85cm] (du) at (0,2.0) {DU};
    \node[mbox,minimum width=3.0cm,draw=ranblue,fill=ranblue] (mb) at (0,0.3)
      {RANBooster DAS middlebox};
    \node[dunode,minimum width=1.4cm,minimum height=0.8cm] (ru1) at (-2.4,-1.7) {RU};
    \node[dunode,minimum width=1.4cm,minimum height=0.8cm] (ru2) at (2.4,-1.7) {RU};
    \draw[wire] (du) -- (mb);
    \draw[wire] (mb.south) -- (ru1.north);
    \draw[wire] (mb.south) -- (ru2.north);

    % Downlink (orange, pointing away from the DU): one packet in, two out.
    \node[pkt] at (-0.75,1.2) {P1};
    \draw[-{Latex[length=2mm]},thick,coreorange] (-1.2,1.5) -- (-1.2,0.9);
    \node[pkt] at (-3.0,-0.75) {P1};
    \node[pktb] at (-2.4,-0.75) {P2};
    \draw[-{Latex[length=2mm]},thick,coreorange] (-3.5,-0.45) -- (-3.5,-1.05);
    \node[pkt] at (2.2,-0.75) {P1};
    \node[pktb] at (2.8,-0.75) {P3};
    \draw[-{Latex[length=2mm]},thick,coreorange] (1.7,-0.45) -- (1.7,-1.05);

    % Uplink (green, pointing back at the DU): two packets in, one combined out.
    \node[pktb,minimum width=1.05cm] at (0.9,1.2) {P2 + P3};
    \draw[-{Latex[length=2mm]},thick,softgreen] (1.65,0.9) -- (1.65,1.5);
    \draw[-{Latex[length=2mm]},thick,softgreen] (-1.85,-1.05) -- (-1.85,-0.45);
    \draw[-{Latex[length=2mm]},thick,softgreen] (3.35,-1.05) -- (3.35,-0.45);

    \node[font=\scriptsize\bfseries,text=ranblue,align=center] at (-3.7,0.7)
      {A2: packet\\replication};
    \node[font=\scriptsize\bfseries,text=ranblue,align=center] at (3.8,0.7)
      {A3 + A4: caching\\ \& modification};
    \node[font=\scriptsize,align=left,text=softgray] at (3.1,2.1)
      {$\sim$850 lines of DPDK code\\$\sim$1300 lines of XDP code};
  \end{tikzpicture}

  \vspace{0.5em}
  \only<2->{\takeaway{The same binary, with no code changes, ran over the SD-RAN,\\
    Capgemini, Radisys and Foxconn stacks --- the fronthaul is standard.}}
\end{frame}

% Teaching note (180 s): run this as a decision, not a result. Give the students the
% constraint (4 radios, 100 MHz, one floor) and make them pick A or B before you
% reveal C. Both A and B are defensible; that is why the heatmaps are convincing.
\begin{frame}[t]{RANBooster DAS benefits: cell planning}
  \vspace{0.1em}
  {\small\textbf{Microsoft Cambridge 5G testbed --- one floor, 4 RUs, 100 MHz of
    spectrum}\par}

  \vspace{0.4em}
  {\small\textbf{How can I do the cell planning?}\par}
  \vspace{0.2em}
  {\small\color{talkred}\textbf{A:} Split 100 MHz into 4 $\times$ 25 MHz cells ---
    no interference, but peak throughput is capped\par}
  \only<2->{{\small\color{ranblue}\textbf{B:} Reuse 100 MHz in all 4 cells ---
    full bandwidth, but the cells interfere\par}}
  \only<3->{{\small\color{softgreen}\textbf{C:} RANBooster DAS: one 100 MHz cell,
    4 radios acting as its antennas\par}}

  \vspace{0.9em}
  \centering
  \begin{tikzpicture}[scale=1.08,transform shape]
    % Panel A: four 25 MHz cells. No interference anywhere, but every cell is capped.
    \begin{scope}[shift={(0,0)}]
      \begin{scope}
        \clip (0,0) rectangle (3.0,1.6);
        \fill[gray!8] (0,0) rectangle (3.0,1.6);
        \foreach \x/\y in {0.5/1.15, 2.5/1.15, 0.5/0.45, 2.5/0.45}{
          \fill[coreorange,opacity=0.26] (\x,\y) circle (0.62);}
      \end{scope}
      \node[font=\small\bfseries,text=talkred] at (-0.28,0.8) {A};
    \end{scope}
    % Panel B: full bandwidth everywhere, so the cells collide.
    \only<2->{
      \begin{scope}[shift={(3.55,0)}]
        \begin{scope}
          \clip (0,0) rectangle (3.0,1.6);
          \fill[gray!8] (0,0) rectangle (3.0,1.6);
          \foreach \x/\y in {0.5/1.15, 2.5/0.45}{
            \fill[softgreen,opacity=0.45] (\x,\y) circle (0.60);}
          \foreach \x/\y in {2.5/1.15, 0.5/0.45}{
            \fill[talkred,opacity=0.24] (\x,\y) circle (0.60);}
        \end{scope}
        \node[font=\small\bfseries,text=ranblue] at (-0.28,0.8) {B};
      \end{scope}}
    % Panel C: one cell, four antennas, uniform across the floor.
    \only<3->{
      \begin{scope}[shift={(7.1,0)}]
        \fill[softgreen,opacity=0.34] (0,0) rectangle (3.0,1.6);
        \node[font=\small\bfseries,text=softgreen] at (-0.28,0.8) {C};
      \end{scope}}
    % Floor outline and the four radios, drawn on top of each panel.
    \foreach \dx/\show in {0/1, 3.55/2, 7.1/3}{
      \begin{scope}[shift={(\dx,0)}]
        \draw[softgray,visible on=<\show->] (0,0) rectangle (3.0,1.6);
        \draw[softgray!50,visible on=<\show->] (1.0,0) -- (1.0,1.6);
        \draw[softgray!50,visible on=<\show->] (2.0,0) -- (2.0,1.6);
        \draw[softgray!50,visible on=<\show->] (0,0.8) -- (3.0,0.8);
        \foreach \x/\y in {0.5/1.15, 2.5/1.15, 0.5/0.45, 2.5/0.45}{
          \draw[talkred,thin,visible on=<\show->]
            (\x-0.09,\y-0.09) -- (\x,\y+0.13) -- (\x+0.09,\y-0.09);}
      \end{scope}}
  \end{tikzpicture}

  \vspace{0.3em}
  {\scriptsize\color{softgray}Downlink throughput across the floor, redrawn
   schematically from the measured heatmaps in the talk.\par}
\end{frame}
